% 强连通分量 = 前向可达集 ∩ 后向可达集。弧：1->2, 2->1, 2->3, 3->4, 4->3, 5->6, 6->4, 6->7, 7->1, 7->5。
% 从 1 前向可达：{1,2,3,4}；能到达 1：{1,2,5,6,7}；交集 {1,2}。
\begin{tikzpicture}[panelmark/.style={font=\figurefont\fontsize{8}{10}\selectfont,text=FigureInk,anchor=north,inner sep=0pt,outer sep=0pt},
font=\figurefont\fontsize{8}{10}\selectfont,>=Stealth,
  vtx/.style={circle,draw=FigureStroke,line width=.8pt,minimum size=6.2mm,inner sep=0pt,text=FigureInk,fill=white},
  arc/.style={draw=FigureStroke,line width=.8pt,->,shorten >=.5pt},
  root/.style={double,double distance=.9pt},
  note/.style={font=\FigureNoteFont\fontsize{7.5}{9}\selectfont,text=FigureMuted,align=center}]
\def\sccgraph#1#2#3{%
  \node[vtx,fill=#1] (v5) at (0,0) {5};
  \node[vtx,fill=#1] (v7) at (0.95,0.95) {7};
  \node[vtx,fill=#1] (v6) at (0.95,-0.95) {6};
  \node[vtx,root,fill=#2] (v1) at (2.35,0.95) {1};
  \node[vtx,fill=#2] (v2) at (3.55,0.95) {2};
  \node[vtx,fill=#3] (v4) at (2.35,-0.95) {4};
  \node[vtx,fill=#3] (v3) at (3.55,-0.95) {3};
  \draw[arc] (v5)--(v6); \draw[arc] (v6)--(v7); \draw[arc] (v7)--(v5);
  \draw[arc,bend left=18] (v1) to (v2); \draw[arc,bend left=18] (v2) to (v1);
  \draw[arc,bend left=18] (v3) to (v4); \draw[arc,bend left=18] (v4) to (v3);
  \draw[arc] (v7)--(v1); \draw[arc] (v2)--(v3); \draw[arc] (v6)--(v4);}
\begin{scope}[local bounding box=panelrow]
\begin{scope}[local bounding box=panela]
  \sccgraph{white}{FigureGreenFill}{FigureGreenFill}

\end{scope}
\begin{scope}[xshift=4.75cm,local bounding box=panelb]
  \sccgraph{FigureBlueFill}{FigureBlueFill}{white}

\end{scope}
\begin{scope}[xshift=9.5cm,local bounding box=panelc]
  \sccgraph{white}{FigureYellowFill}{white}

\end{scope}
\end{scope}
% Subfigure marks: 8 pt, centred below each panel; common top 3 mm below this row.
\coordinate (panelrowbase) at ($(panelrow.south)+(0,-3mm)$);
\node[panelmark] at (panela.center |- panelrowbase) {(a)};
\node[panelmark] at (panelb.center |- panelrowbase) {(b)};
\node[panelmark] at (panelc.center |- panelrowbase) {(c)};
\end{tikzpicture}
